\chapter*{Licence et marque}
\addcontentsline{toc}{chapter}{Licence et marque}

\section*{Ce qui est sous licence LHUNA}
Le nom \emph{LHUNA}, le logo, l'identité visuelle et le présent manuel sont la propriété de leur auteur et sont
diffusés sous licence LHUNA. Toute reproduction du logo ou de ce manuel demande son accord.

\section*{Ce qui est sous licence libre}
L'application \soft{} est un logiciel libre distribué sous \textbf{GNU General Public License, version 3}. Ce choix
n'est pas arbitraire : il découle des composants qu'elle contient.
\begin{itemize}
  \item \soft{} est dérivé de \emph{ZXTune}, distribué sous GNU Lesser General Public License version~3.
        La LGPL~v3 n'autorise une combinaison qu'avec la GPL~v3, pas avec la GPL~v2.
  \item Plusieurs moteurs d'émulation intégrés sont sous GPL~v3 (cœurs Saturn et PlayStation
        \emph{Highly Theoretical/Experimental}) ou GPL~v2 « ou ultérieure » (libsidplayfp, z80ex, ASAP), compatibles
        avec la GPL~v3.
  \item Les autres composants sont sous licences permissives (MIT, BSD, ISC, zlib, MPL~2.0…), toutes compatibles.
\end{itemize}
Conséquence : quiconque reçoit l'APK a le droit d'obtenir le code source correspondant, de le modifier et de le
redistribuer sous les mêmes conditions. La liste complète des composants et de leurs licences figure en
annexe~\ref{chap:credits}.

\section*{Les drivers propriétaires}
Les drivers musicaux de l'auteur (famille \emph{NebulaWeave}) ne font pas partie de l'APK de base. Ils sont
distribués dans un APK séparé, sous licence propriétaire.
\begin{itemize}
  \item Sans cet APK, \soft{} fonctionne exactement comme décrit dans ce manuel.
  \item Avec cet APK installé, \soft{} peut utiliser les drivers qu'il fournit.
\end{itemize}

\begin{technique}
Pour que les deux licences cohabitent, l'APK des drivers doit rester une œuvre séparée au sens de la GPL~v3 :
il fonctionne dans son propre processus et échange avec \soft{} par une interface documentée (service Android,
fournisseur de contenu, données). Charger sa bibliothèque native dans le processus de \soft{} ferait de l'ensemble une
seule œuvre, qui devrait alors être entièrement sous GPL~v3.
\end{technique}

\section*{Marques citées}
Les noms de consoles, d'ordinateurs, de puces et de jeux cités dans ce manuel sont la propriété de leurs détenteurs
respectifs et ne sont utilisés que pour décrire la compatibilité.
